\documentclass[12pt]{article}

\addtolength{\textwidth}{.8in}
\addtolength{\hoffset}{-.4in}
\addtolength{\textheight}{1in}
\addtolength{\voffset}{-1in}

\usepackage{amsmath}
\usepackage{graphicx}
\usepackage{upquote} % to convert funny quotes to straight quotes
\usepackage{amsfonts}
\usepackage{paralist}
\setlength{\parindent}{15pt}

\renewcommand*{\thefootnote}{\fnsymbol{footnote}}

% Plan for class:

% 1. LOESS w/ students at board doing best-fit lines for each of 10 data
% points, (alpha = .4, uniform kernel).
% Switch to best-fit *constant* (not line) and note that this is moving
% average.
%
% 2. Merging: on the board, movies/actors example from notes.
% pd.merge(df1, df2, left_on='', right_on='')
% Also do Halloween example (many-to-many)
%
% On screen: Pokemon data set. "What colors of Pokemon can do karate chops?"
%
% 3. Tidy data: on board, define: variable, value, observation, observational
% unit. (p.6) Then do this example:
% 
% Name     Gen drunk sober
% Jezebel   F   80   90
% Biff      M   60   60
% Filbert   M   85  110
% Betty Lou F  135  160
%
% Break into:
%
% Name    state  IQ                      Name     Gen
% Jezebel drunk  80                      Jezebel   F
% ...                                       ...
%
% tdf = pd.melt(df, ['Name'], var_name='state', value_name='IQ')   "long form"
% to put in "wide form":
% df = tdf.pivot(index='Name', columns='state')
%
% Exercise on last page.

\begin{document}

\begin{center}
\begin{Large}
DATA 219 Notes -- Feb. 13, 2018
\end{Large}
\end{center}

\section{Merge operations (and data fusion)}

As we've learned, the \textit{lingua franca} of the data science world is the
rectangular table perfectly reflected by a CSV file: it has (heterogeneous)
rows, it has (homogeneous) columns, and is otherwise flat. Nearly all machine
learning algorithms (including those for prediction and modeling) are based on
such a structure. In Pandas we speak of this structure as a DataFrame, of
course, and we'll sometimes also use the word \textbf{table} to describe this
conceptual format for organizing data.

More complicated data sets sometimes consist of \textit{multiple} such tables.
This is especially true if there are different aspects to the data, each of
which is deserving of its own table. The tables are often implicitly related,
as we'll see, through some piece of common information that links them
together logically.

Moreover, there will be times in your career when you want to integrate
multiple \textit{different} data sets in some way, perhaps because each gives
you information on some related aspect of a common phenomenon. Imagine if you
found a data source giving vote totals for various political candidates, and a
different data source documenting the financial donations that their campaigns
received. An obvious question would be: ``to what extent do a candidate's
contributions relate to his/her vote count?" This kind of analysis is only
possible if the two data sets can be merged in a coherent way, so that the
vote totals for a candidate are properly associated with the contributions to
that same candidate's campaign.

We'll use the term \textbf{merge} (or \textbf{join}) to describe the operation
of connecting two tables/DataFrames into one, based on a common linking piece
of information. The term \textbf{fusion} is usually restricted to the case
described in the previous paragraph; namely, when the two tables to be joined
are from different data sets entirely.

\subsection{Merging in Pandas}

The basic idea is that in order to meaningfully merge two DataFrames, one (or
more) columns in the first table must be logically linked to one (or more)
columns in the second. ``Logically linked" usually means ``having the same
values in them." For example, consider these two DataFrames from a data source
like IMDB:

\pagebreak
\begin{Verbatim}[fontsize=\small,samepage=true,frame=none]
                        actors:
actor           film                   salary (millions)
--------------------------------------------------------
Downey Jr.      The Avengers           50
Johansson       The Avengers           20
DiCaprio        Inception              59
Hanks           Forrest Gump           70
Bullock         Gravity                77
Downey Jr.      Iron Man 3             51
Hanks           Saving Private Ryan    40
...

                        movies:
title           director               year
-------------------------------------------
Inception       Nolan                  2010
The Avengers    Whedon                 2012
Forrest Gump    Zemeckis               1994
Gravity         Cuaron                 2013
Star Wars       Lucas                  1977
...
\end{Verbatim}
\normalsize
 
There are clearly questions we'd like to ask which involve both DataFrames
together, rather than each in isolation. For instance, ``have actor salaries
been going up over time?" or ``how often have various actors worked with
specific directors?"

Merging here is simple: we simply identify \textit{the common column} which
unites rows from the two tables. Clearly, it's the film title: we know that
Robert Downey Jr. was directed by Joss Whedon because a Downey-Jr. row in the
\texttt{actors} table has the same film name as a film in a Whedon row in the
\texttt{movies} table.

In database parlance, we call the \texttt{title} column of the \texttt{movies}
table its \textbf{primary key} (PK), and we call the \texttt{film} column of
the \texttt{actors} table a \textbf{foreign key} (FK). The primary key /
foreign key relationship is supposed to suggest that for every row in the
\texttt{actors} table, its \texttt{film} value \textit{will} appear in one row
of the \texttt{movies} table (namely, in that row's \texttt{title} column.)

To ``join" these two DataFrames all in one fell swoop, we execute a command
like this:

\begin{Verbatim}[fontsize=\small,samepage=true,frame=none]
actors_movies = pd.merge(actors,movies,left_on="film",right_on="title")
\end{Verbatim}

This produces a \textit{new} DataFrame that contains (1) all the columns of
the two, and (2) a row for each pair of original rows that match on their
movie title. Notice we had to specify \textit{two} arguments (besides the
table names) to \texttt{pd.merge()}: ``\texttt{left\_on}" (for the name of the
merging column in the first DataFrame) and ``\texttt{right\_on}" (for the
second). If these two columns had the same name, we could just specify a
single argument called ``\texttt{on}" for the two.

\vspace{.2in}
Here's what \texttt{actors\_movies} will look like:
\begin{Verbatim}[fontsize=\small,samepage=true,frame=none]
                            actors_movies:
actor      film       salary (millions)   title      director     year
----------------------------------------------------------------------
Bullock    Gravity    77                  Gravity    Cuaron       2013
DiCaprio   Inception  59                  Inception  Nolan        2010
....
\end{Verbatim}
(The sequence of the rows is not in general reliable, and you should not
depend on them being in the order you expect.)

This DataFrame can be used to answer either of the questions we posed above,
and others.

Note that our new table has some redundant information (the \texttt{film} and
\texttt{title} columns are identical) which you can remove using normal Pandas
operations if you wish.

\subsubsection{Odds and ends}

\begin{enumerate}
\itemsep.1em
\item Sometimes we want to merge on more than one column; for instance:

\begin{Verbatim}[fontsize=\small,samepage=true,frame=none]
df = pd.merge(votes,contributions,left_on=['first_name','last_name'],
        right_on=['candidate_first','candidate_last'])
\end{Verbatim}
which has the expected result.

\vspace{.2in}
\item The default merge operation (as we did above) is sometimes called an
``\textbf{inner join}". This means that \textit{if there is no corresponding
match for a row in one of the DataFrames, there will be no associated row in
the result.} This is normally what we want. In the movies example, if we don't
have a row in the \texttt{movies} table for a particular film, then we don't
know the year for that film, and it's (perhaps) pointless to include any of
the actors in that film in the resulting merged DataFrame.

Sometimes, we want to make sure that every row in the left (or right, or
both) DataFrame appears in the result \textit{at least once}, though, with
corresponding null values if necessary. We can do this by specifying
"\texttt{left}", "\texttt{right}", or "\texttt{outer}" as the value of the
``\texttt{how}" argument to \texttt{pd.merge()}:

\begin{Verbatim}[fontsize=\small,samepage=true,frame=none]
df = pd.merge(votes,contributions,left_on=['first_name','last_name'],
        right_on=['candidate_first','candidate_last'],how='outer')
\end{Verbatim}

\vspace{.2in}
\item Sometimes merging two data frames will result in column name
collisions. Suppose, for example, that our \texttt{actors} table had had a
\texttt{year} column indicating the birth year of the actor:

\begin{Verbatim}[fontsize=\small,samepage=true,frame=none]
                        actors:
actor           year   film                   salary (millions)
---------------------------------------------------------------
Downey Jr.      1965   The Avengers           50
Johansson       1984   The Avengers           20
DiCaprio        1974   Inception              59
Hanks           1956   Forrest Gump           70
Bullock         1964   Gravity                77
...
\end{Verbatim}

Then the resulting DataFrame would need to have two columns named
``\texttt{year}", with different meanings. Pandas solves this by renaming the
conflicting columns with a ``\texttt{\_x}" or ``\texttt{\_y}" suffix,
corresponding to the left/right tables in the join:

\footnotesize
\begin{Verbatim}[fontsize=\small,samepage=true,frame=none]
                            actors_movies:
actor      year_x   film       salary (millions)   title      director   year_y
-------------------------------------------------------------------------------
Bullock    1964     Gravity    77                  Gravity    Cuaron     2013
DiCaprio   1974     Inception  59                  Inception  Nolan      2010
\end{Verbatim}
\normalsize

(which, of course, you can rename at will via ``\texttt{actors\_movies.columns}".)

\item Finally, the above examples have all featured ``one-to-many" relationships
between the two tables. (\textit{Many} actors star in \textit{one} movie.) But
sometimes we have a ``many-to-many" relationship. Consider the following
example:

\small
\begin{Verbatim}[fontsize=\small,samepage=true,frame=none]
costumes = pd.DataFrame(
    {'costume':['Darth Vader','Darth Vader','Princess','Princess'],
     'accessory':['helmet','lightsaber','crown','wand']})

    accessory      costume
0      helmet  Darth Vader
1  lightsaber  Darth Vader
2       crown     Princess
3        wand     Princess

children = pd.DataFrame(
    {'child':['Johnny','Sam','Betty Lou','Wendell','Kiara'],
     'costume':['Darth Vader','Darth Vader','Princess','Student','Princess']})

       child      costume
0     Johnny  Darth Vader
1        Sam  Darth Vader
2  Betty Lou     Princess
3    Wendell      Student
4      Kiara     Princess

pd.merge(costumes,children,on="costume")

    accessory      costume      child
0      helmet  Darth Vader     Johnny
1      helmet  Darth Vader        Sam
2  lightsaber  Darth Vader     Johnny
3  lightsaber  Darth Vader        Sam
4       crown     Princess  Betty Lou
5       crown     Princess      Kiara
6        wand     Princess  Betty Lou
7        wand     Princess      Kiara
\end{Verbatim}
\normalsize

In order to dress every child in their Halloween costume, we need
\textit{multiple} matches from \textit{both} data frames for each costume
value. Each \texttt{Princess} has more than one accessory, and there are
multiple children who want to \textit{be} princesses. Therefore, we can
potentially end up with a lot of rows in our result.

Btw, notice that poor \texttt{Wendell} was left out of the merged DataFrame
entirely, since his costume did not have any accessories. We can include him
by doing a ``right join" (or full ``outer join"):

\begin{Verbatim}[fontsize=\small,samepage=true,frame=none]
pd.merge(costumes,children,on="costume",how="outer")

    accessory      costume      child
0      helmet  Darth Vader     Johnny
1      helmet  Darth Vader        Sam
2  lightsaber  Darth Vader     Johnny
3  lightsaber  Darth Vader        Sam
4       crown     Princess  Betty Lou
5       crown     Princess      Kiara
6        wand     Princess  Betty Lou
7        wand     Princess      Kiara
8         NaN      Student    Wendell
\end{Verbatim}

\end{enumerate}

\pagebreak
\subsection{``By the way...why?"}

At this point you may be asking, ``if we just have to merge the DataFrames to
do our analysis, why were they stored that way in the first place? Why not
just have everything in one big CSV file?"

The answer has to do with \textit{redundancy of information}. Consider the
merged \texttt{actors\_movies} table. It's not the most compact way of
expressing the information, because certain pieces of information are repeated
unnecessarily.

For instance, \textit{The Avengers} was released in 2012. Our data source
needs to record that. But in the merged table, this year would be repeated
\textit{for every actor/actress in the movie.} The same is true for the
director of the film. By performing the merge, we have duplicated lots of
values. In database terms, the data set has become \textbf{unnormalized}. This
isn't always a bad thing, and in terms of performing standard ML algorithms,
it's necessary. But you can see why we wouldn't want the data permanently
stored that way. It consumes extra storage, plus it opens a Pandora's Box of
possible contradictions: what if \texttt{Downey Jr.}'s \texttt{The Avengers}
row said 2012, while \texttt{Johansson}'s said 2014. Who would we believe?

Database normalization is a larger topic than we have time to fully cover in
this class, but the basics are very simple. Conceptually: \textit{each piece
of information should be stored in one and only one place.} Generally
speaking, when that's true, you have a \textbf{normalized} data set, which is
most appropriate for storage and for communication to other potential users.
Analysts will proceed to unnormalize select parts of it in order to perform
certain analyses.

\section{``Tidy" data}

The term ``tidy" to describe a data set has sprung into prominence recently as
a result of the excellent paper by Hadley Wickham (of R fame) who neatly
captured the aspects of various data formatting choices that make life
problematic for the analyst. 

\begin{compactitem}
\item \textbf{variable} -- a category of something, for which individual
objects have a value. For example: height, gender, center\_field\_distance,
name.
\item \textbf{value} -- an actual individual piece of data corresponding to a
variable for one object. For example: 5'6", female, 390 feet, Kylo Ren.
\item \textbf{observation} -- a particular measurement that comprises all the
information necessary to interpret that measurement. For example: Jeff's
weight was 165 lbs., Jenny's golf score at Field \& Cross Links in the rain
was 88. (In the first example, ``all the information necessary" includes
``Jeff" and ``165"; in the second, it includes ``Jenny", ``Field \& Cross
Links", ``rainy", and ``88".)
\item \textbf{observational unit} -- a collection of all the observations of a
particular type. (``Golf performances" would be an observational unit,
separate from a ``golfers" unit which contains information about golfers
themselves, not their particular performances.)
\end{compactitem}

\vspace{.2in}
The rule is: a data set is considered \textbf{tidy} if and only if:
\begin{compactenum}
\item Each observational unit is in its own table. (This is essentially
the same as ``normalized," from the previous section.)
\item In each table, each observation forms a row.
\item In each table, each column header is a variable, and the column contains
values.
\end{compactenum}

\vspace{.2in}
This is best visualized by looking at examples of \textit{un}tidy data sets
and rectifying them. Here's one:
\small
\begin{Verbatim}[fontsize=\small,samepage=true,frame=none]
            religion   <$10k   $10-20k  $20-30k  $30-40k   $40-50k  $50-75k
0           Agnostic      27        34       60       81        76      137
1            Atheist      12        27       37       52        35       70
2           Buddhist      27        21       30       34        33       58
3           Catholic     418       617      732      670       638     1116
4  Dont know/refused      15        14       15       11        10       35
...
\end{Verbatim}
\normalsize

The problem here is that \textit{the column headings are actually values, not
variables}. Think about it: ``\texttt{\$10-20k}" isn't the name of a variable
which might have actual values, but the value (range) itself of some variable.

Two questions: (1) What would the variable be for these column values? Answer:
\textit{Income.} (2) What are the cells themselves representing? Answer:
\textit{Frequency.} (\textit{i.e.}, the number of people polled in the survey
who had that religion and that income level.)

We tidy this with the \texttt{pd.melt()} function:

\begin{Verbatim}[fontsize=\small,samepage=true,frame=none]
tdf = pd.melt(df, ['religion'], var_name='income', value_name='frequency')
\end{Verbatim}

which gives us the tidy version:
\begin{Verbatim}[fontsize=\small,samepage=true,frame=none]
                    religion    income  frequency
0                   Agnostic     <$10k         27
1                    Atheist     <$10k         12
2                   Buddhist     <$10k         27
3                   Catholic     <$10k        418
    ...
\end{Verbatim}
(Sometimes you'll hear the terms ``wide" and ``long" to refer to the untidy
and tidy formats in this sort of example, respectively.)

\vspace{.1in}
By the way, the inverse operation of ``melt" is often called
``\textbf{pivot}". To reverse the above transformation in Pandas, you should
use:

\begin{Verbatim}[fontsize=\small,samepage=true,frame=none]
df = tdf.pivot(index='religion',columns='income')
\end{Verbatim}

\vspace{.1in}
\subsection{Exercise}

Test your understanding of tidy datasets by putting the following (messy) set
into tidy form:

\small
\begin{Verbatim}[fontsize=\small,samepage=true,frame=none]
Name  Year  Major  School  Type    GPAFr  GPASo  GPAJr  GPASr  Creds  Mascot
----------------------------------------------------------------------------
Jane  Soph   CPSC     UMW  Public    2.4    3.8     NA     NA     48   Eagle
Fred  Sen    ENGL     JMU  Public    2.9    3.1    2.3    2.9    106   Duke 
                                  ...
\end{Verbatim}
\normalsize

% Exercise:
% load http://cs.umw.edu/~stephen/data219/students.csv and tidy it to:
%
%                      students
%            name  year major    school  creds
%   Mal A. Kenobi  Soph  ECON       UMW     32
%  Jay O. Aniston   Sen  ECON  Richmond     92
%
%               gpas
%             name year   GPA
%   Mal A. Kenobi   So  2.47
%  Jay O. Aniston   Sr  2.30
%      (no NaNs)
%
%           schools
%    school     type   mascot
%       UMW   public   Eagles
%  Richmond  private  Spiders
%
% DELETE the original table.
% now add a column for "urban/rural" to schools table. (We'll call Richmond,
% Fredericksburg, and Charlottesville urban, and the others rural.)
%
% how many students are there?
% how many total year-GPA measurements are there?
% what's the overall GPA?
% how many students in rural vs. urban schools?
% average freshman GPA in rural vs. urban schools?
% visual: do freshmen vs. non-freshmen GPAs differ? (will require another col)
% visual: does avg GPA vary by major? by school? (will require another col)
% scatterplot of creds vs. average-yearly-GPA (ditto)
%
% helpful:
% 1) after splitting out school into separate df:
%  df.drop_duplicates()
% 2) to remove NaN's from GPAs once melted:
%  np.isnan(...)   <-- a ufunc
%  ~   <-- boolean "not" operator
%  (to set up: "gpas = gpas[~np.isnan(gpas.GPA)]")
%  (could also do: "gpas.drop(gpas.index[np.isnan(gpas.GPA)])")
% 3) df.col.str   to perform string operations as ufunc
%  (to set up: "gpa.year = gpa.year.str[3:]")
%
% 


